\subsection{4. Network Cards}

\textbf{Dispositivos físicos --- host principal (Windows):}

\begin{table}[H]
  \centering
  \begin{tabular}{ll}
    \toprule
    Campo & Valor \\
    \midrule
    Adaptador Wi-Fi & Realtek 8852CE WiFi 6E PCI-E NIC \\
    MAC Wi-Fi & 08-F9-7E-9E-25-03 \\
    IPv4 & 192.168.1.10/24 (DHCP) \\
    Gateway / DNS & 192.168.1.1 \\
    SSID / enlace & EVANGELIO\_5G (5 GHz) --- 866.7 Mbps \\
    Bytes RX / TX & 3.250.114.444 / 3.535.449.942 \\
    Adaptador Ethernet & Realtek PCIe GbE Family Controller (sin cable) \\
    \bottomrule
  \end{tabular}
  \caption{Ficha de la tarjeta de red del host.}
\end{table}

Otros adaptadores virtuales del host: Tailscale Tunnel, VirtualBox Host-Only (192.168.56.1) y 2 adaptadores Wi-Fi Direct.

\textbf{Máquinas virtuales (2 comparadas con el host):}

\begin{table}[H]
  \centering
  \begin{tabular}{lll}
    \toprule
    Campo & Slackware 15.0 & Solaris 11.4 \\
    \midrule
    NIC & VirtualBox e1000 & e1000g0 (Intel PRO/1000) \\
    MAC & 08:00:27:0a:46:19 & virtual (dladm show-phys) \\
    IPv4 (casa) & 192.168.1.13 (DHCP) & 192.168.1.11 (DHCP) \\
    Velocidad & virtual & 1000 Mb/s full-duplex \\
    \bottomrule
  \end{tabular}
  \caption{Datos de red de las VMs comparadas.}
\end{table}

\textbf{Comparación host vs VMs:} el host usa NICs físicas Realtek (Wi-Fi 6E + GbE); las VMs usan NICs virtuales emuladas. Las MAC de las VMs inician con \texttt{08:00:27} (OUI de VirtualBox), lo que demuestra que son virtuales. La velocidad del enlace Solaris es 1000 Mbps full-duplex (virtual), equivalente a la GbE física del host.

\subsection{5. Shell Programming --- Unix (Slackware 15.0)}

Se escribieron y ejecutaron los 4 scripts pedidos en Slackware 15.0 (guardados en \texttt{/root/scripts/}). Todos fueron probados con salidas reales el 17/08/2026.

\subsubsection{Script 1.1 --- mi\_ls.sh (menú de listado)}

Lista los archivos de un directorio con menú interactivo: orden por fecha/tamaño, conteo por grupos, filtros por nombre y recursión.

Ejecución real --- opción 5 (contar por tipo) sobre \texttt{/etc}:

\begin{lstlisting}
======================================
  MENU - Listar archivos de: /etc
======================================
1) Mas recientes (conteo por fecha)
2) Mas antiguos (conteo por fecha)
3) Mayor tamano (conteo por tamano)
4) Menor tamano (conteo por tamano)
5) Por tipo (archivo/directorio)
6) Filtro: empieza con...
7) Filtro: termina con...
8) Filtro: contiene...
9) Incluir subdirectorios (recursivo)
0) Salir
======================================
Directorio: /etc
--- Directorios (con ocultos):
112
--- Archivos (con ocultos):
172
\end{lstlisting}

Resultado: \texttt{/etc} tiene 112 directorios y 172 archivos (incluyendo ocultos).

Ejecución real --- opción 6 (filtro: empieza con ``rc''):

\begin{lstlisting}
/etc/rc_keymaps
/etc/rc3.d
/etc/rc5.d
/etc/rc0.d
/etc/rc1.d
/etc/rc4.d
/etc/rc6.d
/etc/rc_maps.cfg
/etc/rc2.d
/etc/rc.d
\end{lstlisting}

Comandos clave: \texttt{ls -lat / -latr} (por fecha), \texttt{ls -laS / -laSr} (por tamaño), \texttt{awk} (extraer columnas de fecha), \texttt{sort | uniq -c} (agrupar y contar), \texttt{grep "\^{}d"} (detectar directorios), \texttt{find -name} con \texttt{-maxdepth 1}, \texttt{wc -l}.

\subsubsection{Script 1.2 --- buscar.sh (búsqueda)}

Menú de búsqueda: por nombre de archivo, por palabra dentro de archivo, combinado, conteo de líneas, head/tail.

Ejecución real --- buscar la palabra ``dev'' dentro de \texttt{/etc/fstab}:

\begin{lstlisting}
--- Coincidencias (linea: contenido):
1:/dev/sda3   /         ext4   defaults   1  1
2:/dev/sda1   /boot     ext4   defaults   1  2
3:/dev/sda2   swap      swap   defaults   0  0
--- Total de ocurrencias:
\end{lstlisting}

Resultado: encontró las 3 líneas de fstab que mencionan \texttt{/dev/}, con número de línea (gracias a \texttt{grep -n}).

Comandos clave: \texttt{find dir -name "*patrón*"}, \texttt{grep -n}, \texttt{grep -c}, \texttt{wc -l}, \texttt{head/tail -n N}, \texttt{2>/dev/null} (silenciar errores de permiso).

\subsubsection{Script 1.3 --- revisar\_logs.sh (logs)}

Muestra las últimas 15 líneas de \texttt{/var/log/syslog}, \texttt{/var/log/messages} y \texttt{/var/log/secure}, con filtro opcional.

Verificación de los tres logs:

\begin{lstlisting}
-rw-r----- 1 root root 1255029 Aug 17 21:13 /var/log/messages
-rw-r----- 1 root root  290449 Aug 17 21:13 /var/log/secure
-rw-r----- 1 root root  402441 Aug 17 21:13 /var/log/syslog
\end{lstlisting}

Ejecución real --- filtro por ``sshd'':

\begin{lstlisting}
===== /var/log/messages (filtro: sshd) =====
Aug 17 21:13:02 slackware sshd[2007]: Received disconnect from 10.0.2.2 port 55807:11: disconnected by user
Aug 17 21:13:02 slackware sshd[2007]: Disconnected from user vagrant 10.0.2.2 port 55807
Aug 17 21:13:12 slackware sshd[2025]: Accepted publickey for vagrant from 10.0.2.2 port 62151 ssh2: RSA SHA256:...
\end{lstlisting}

\textbf{Respuestas sobre logs:}
\begin{enumerate}
  \item \textbf{¿Qué son los archivos de log?} Registros donde el SO y los servicios guardan eventos (fecha, origen, mensaje). Sirven para diagnóstico, auditoría y seguridad.
  \item \textbf{¿Qué tipos de logs hay en los SO instalados?} Slackware: \texttt{syslog}, \texttt{messages}, \texttt{secure} (+ dmesg/kernel). Solaris: \texttt{/var/adm/messages}. Windows: Visor de eventos (System, Security, Application).
  \item \textbf{¿Qué es syslog y qué define el estándar?} Protocolo/estándar de logging (RFC 5424): define formato (facilidad + severidad), transporte (UDP 514) y el servicio; centraliza logs de distintos dispositivos.
  \item \textbf{¿Los logs encontrados siguen el estándar?} Sí: Slackware/Solaris usan formato syslog (facilidad.severidad, timestamp, host, proceso[pid]: mensaje). Windows usa Event Log propio (no syslog nativo).
\end{enumerate}

\subsubsection{Script 1.4 --- newgroup.sh y newuser.sh}

\texttt{newgroup.sh}: crea un grupo con GID explícito validando que no exista.

\begin{lstlisting}
$ sudo bash newgroup.sh ventas 1099
Grupo ventas creado (GID 1099)
$ getent group ventas
ventas:x:1099:
$ sudo bash newgroup.sh ventas 1099
ERROR: el GID 1099 ya existe en el sistema.    <- evita duplicados
\end{lstlisting}

\texttt{newuser.sh}: automatiza usuario + grupo + home + subdirectorios + permisos.

\begin{lstlisting}
$ sudo bash newuser.sh alice developers "Alice Developer" /home/alice /bin/bash 700 770 755
Creando grupo developers...
Usuario alice creado (home: /home/alice, grupo: developers, shell: /bin/bash)
Permisos aplicados: /home/alice=700, documentos=770, scripts=755
Listo.

$ getent passwd alice
alice:x:1005:1004:Alice Developer:/home/alice:/bin/bash

$ sudo ls -la /home/alice/
drwx------ 4 alice developers ...  .          <- 700: solo alice entra
drwxrwx--- 2 root  root      ...  documentos <- 770: grupo puede leer/escribir
drwxr-xr-x 2 root  root      ...  scripts    <- 755: todos pueden leer/ejecutar

$ sudo bash newuser.sh alice developers "Alice" /home/alice /bin/bash 700 770 755
ERROR: el usuario alice ya existe.     <- evita duplicados
\end{lstlisting}

Parámetros del script: \texttt{\$1} usuario, \texttt{\$2} grupo, \texttt{\$3} nombre completo (GECOS), \texttt{\$4} home, \texttt{\$5} shell, \texttt{\$6} permisos home (700), \texttt{\$7} permisos documentos/ (770), \texttt{\$8} permisos scripts/ (755).

\subsection{6. Editor VI --- Ejercicio del Himno}

Archivo de trabajo: \texttt{himno.txt} (16 líneas, 4 estrofas):

\begin{lstlisting}
 1  Estudiante, maestro la conquista
 2  Sera hacer con amor nuestra labor
 3  Cultores de espiritu humanista
 4  Unidad de intelecto y corazon.
 5  Escuela de ingenio es nuestra casa
 6  Libro abierto a nuestra universidad
 7  Aqui perdura mientras todo pasa
 8  Cimiento de la fe y la integridad.
 9  Ofrecemos la mano al que tropieza
10  La hidalguia del dialogo al rival
11  Ofrecemos la duda y la certeza
12  Mediamos entre hierro y el cristal.
13  Escuela de ingenio es nuestra casa
14  Libro abierto a nuestra universidad
15  Aqui perdura mientras todo pasa
16  Cimiento de la fe y la integridad.
\end{lstlisting}

\textbf{Operaciones aplicadas con sus resultados reales:}

\begin{itemize}
  \item \texttt{:1,4s/a/-/g} --- todas las ``a'' de las líneas 1--4 reemplazadas por ``-'': \texttt{1,4} = rango · \texttt{s} = substitute · \texttt{g} = global (todas las de cada línea).
  \item \texttt{:\%s/al/\#\#/g} --- ``al'' por ``\#\#'' en TODO el archivo (reemplazó en 3 líneas); \texttt{\%} = todo el archivo.
  \item \texttt{dw} --- borrar palabra bajo el cursor (borró ``Estudiante'' de la línea 1). \texttt{d}=delete, \texttt{w}=word.
  \item \texttt{:13,16d} --- borrar las líneas 13--16 de una (rango + delete).
  \item \texttt{u} --- deshacer: el archivo volvió a su estado anterior (historial de cambios).
  \item \texttt{G} + \texttt{gUU} --- ir a la última línea y convertirla a MAYÚSCULAS: línea 16 quedó ``CIMIENTO DE LA FE Y LA INTEGRIDAD.''.
  \item \texttt{7G yy yy G p} --- a la línea 7, copiar 2 líneas y pegarlas al final (líneas 16--17 nuevas).
  \item \texttt{/Escuela} --- búsqueda: ``Escuela'' aparece en las líneas 5 y 13; \texttt{n} va a la siguiente coincidencia.
  \item \texttt{5G} --- ir a la línea 5.
  \item \texttt{:wq} --- guardar y salir.
  \item \texttt{:1,5d} --- al reabrir: borrar las primeras 5 líneas.
\end{itemize}

\textbf{Tabla resumen de comandos VI:}

\begin{table}[H]
  \centering
  \begin{tabular}{lll}
    \toprule
    Modo & Comando & Qué hace \\
    \midrule
    Normal & \texttt{i} / \texttt{a} & Insertar antes / después del cursor \\
    Normal & \texttt{Esc} & Volver al modo normal \\
    Normal & \texttt{x} / \texttt{dw} & Borrar carácter / palabra \\
    Normal & \texttt{dd} / \texttt{5dd} & Borrar línea / 5 líneas \\
    Normal & \texttt{u} / \texttt{Ctrl+r} & Deshacer / rehacer \\
    Normal & \texttt{yy} / \texttt{p} & Copiar línea / pegar \\
    Normal & \texttt{G} / \texttt{5G} & Ir al final / a la línea 5 \\
    Normal & \texttt{gUU} & Línea actual a MAYÚSCULAS \\
    Normal & \texttt{/palabra} & Buscar (\texttt{n} siguiente) \\
    Comando & \texttt{:w} & Guardar sin salir \\
    Comando & \texttt{:q} / \texttt{:wq} / \texttt{:q!} & Salir / guardar y salir / salir sin guardar \\
    Comando & \texttt{:\%s/old/new/g} & Reemplazar en todo el archivo \\
    Comando & \texttt{:1,4s/a/-/g} & Reemplazar en rango de líneas \\
    Comando & \texttt{:13,16d} & Borrar rango de líneas \\
    Comando & \texttt{:set number} & Mostrar números de línea \\
    \bottomrule
  \end{tabular}
  \caption{Comandos VI aplicados y verificados en el ejercicio.}
\end{table}
